% !TeX root = writeup.tex 

\section{Conclusion and Future Work}

We argue that without a careful model of delayed outcomes, we
cannot foresee the impact a fairness criterion would have if enforced as a
constraint on a classification system. 
However, if such an accurate outcome model is available, we show that there are more direct 
ways to optimize for positive outcomes than via existing fairness criteria.

Our formal framework exposes a concise, yet expressive way to model outcomes via the expected
change in a variable of interest caused by an institutional decision. This leads to the natural
concept of an outcome curve that allows us to interpret and compare solutions
effectively.
%
In essence, the formalism we propose requires us to understand the two-variable
causal mechanism that translates decisions to outcomes. Depending on the
application, such an understanding might necessitate greater domain knowledge
and additional research into the specifics of the application. This is
consistent with much scholarship that points to the context-sensitive nature of
fairness in machine learning.

An interesting direction for future work is to consider other characteristics of
impact beyond the change in population \emph{mean}. Variance and
individual-level outcomes are natural and important considerations. Moreover, it
would be interesting to understand the robustness of outcome optimization to
modeling and measurement errors.
